%=======================================================================
% SUPERSEDED SECTIONS AND THEIR REPLACEMENTS
%
% Against the draft titled "SAHAi: Socratic RL Tutor" (the one whose
% abstract reports 17.5% on 20 unseen problems).
%
% Each block below is marked:
%   [[R-n]] REPLACE  ... what to delete, named by its opening words
%           WITH     ... the LaTeX to paste in its place
%   [[A-n]] ADD      ... new material, with the insertion point named
%
% No em dashes anywhere, per the house rule.
% Every number here is measured. Provenance is in docs/04-findings.md,
% findings 23 to 27.
%=======================================================================


%-----------------------------------------------------------------------
% [[R-1]] REPLACE: the whole abstract.
%
% Superseded on three counts. The 17.5% headline was measured with a
% verifier that failed 15% of references, an entry point that made 5 of
% 60 held-out problems unwinnable, and an extractor that discarded
% correct code on a name mismatch. The "6% to 18%, no statistical
% difference" framing predates the controls. And "validated tutor policy
% reward design" is not supported: the correctness term did nothing.
%-----------------------------------------------------------------------
\begin{abstract}
A language model that solves programming problems is not thereby able to teach
them. The behaviours conflict: the fastest way to make a student's next
submission correct is to hand over the answer, which is the one thing a tutor
must not do. SAHAi treats this as a reinforcement learning problem. A tutor
policy holds a multi-turn dialogue with a simulated, knowledge-traced student,
the student then attempts the problem, and the attempt is executed against unit
tests. The tutor is rewarded for the student's verified success and penalised
for disclosing the solution. No human tutoring transcripts are used. The design
follows the review of pedagogical reinforcement learning in~\cite{rev-pedrl}.

Our central result is negative, and the most useful part of it is a negative
result about our own measurement. On 60 held-out problems compared pairwise, a
tutored student solved $0.183$ against $0.367$ for the same student given no
dialogue at all (McNemar exact, $p{=}0.0074$), which we previously reported as
evidence that tutoring costs the learner. It does not. Replaying the same
transcripts, with the same tutor turns, while omitting the student's own replies
from the context of its final attempt lifts that arm to $0.350$ and removes the
deficit entirely ($p{=}1.0000$). What survives is weaker and better supported:
tutoring as built neither helps nor harms, and training is inert. The trained
and untrained adapters score identically once the framing is corrected ($0.350$
against $0.350$, $p{=}1.0000$), and the one comparison in which training ever
looked beneficial ($p{=}0.0352$) measured its partial immunity to the artefact
rather than better teaching.

We give the mechanism for the inertia. Every reward term the policy can compute
from its own output is driven to its ceiling, pedagogy $0.76$ to $0.97$ and
leakage $0.44$ to $0.09$, while the outcome term does not move. We also report
four instrumentation defects that depressed every number this project has
published, two of them found by auditing why an oracle arm fails at all: a
verifier that marked $15\%$ of reference solutions wrong through a JSON round
trip, a seed that was never applied, an entry point read from the first function
a reference defines rather than the one its tests call, and an extraction step
that discarded a correct solution whenever the student's function name differed
in case. Correcting the last two raises the oracle ceiling from $0.633$ to
$0.783$ and more than doubles the headroom available to a tutor, from $0.267$ to
$0.417$. Alongside the training results we describe a complete interactive
system: six containerised services, voice input for Indic languages, and a
persistent Bayesian learner model driving placement and curriculum.
\end{abstract}


%-----------------------------------------------------------------------
% [[R-2]] REPLACE: the IEEEkeywords block.
%
% The current block is three sentences of prose, not keywords.
%-----------------------------------------------------------------------
\begin{IEEEkeywords}
Reinforcement learning from verifiable rewards, group-relative policy
optimisation, intelligent tutoring systems, simulated learners, knowledge
tracing, reward hacking, measurement validity, large language models.
\end{IEEEkeywords}


%-----------------------------------------------------------------------
% [[R-3]] REPLACE: in \subsection{Contributions}, items 3 and 4.
%
% Item 3 states the pedagogy term carries about five times the usable
% gradient of the outcome term. That ratio reversed once partial credit
% was added, and keeping it would be reporting a measurement we have
% since overturned. Item 4 is still true but is no longer the headline
% methodological point; the framing artefact is.
%-----------------------------------------------------------------------
\item We measure the asymmetry between reward channels and report that it
      reversed under a change we made to fix it. Under strict all-or-nothing
      scoring the deterministic, tutor-controlled terms carried roughly five
      times the usable between-group variance of the outcome term ($0.0482$
      against $0.0098$). After partial test-case credit the ratio runs the other
      way ($0.0150$ against $0.0481$), and held-out solve rate still does not
      move. Gradient starvation was real and worth fixing. It was not what
      prevented learning (Sec.~\ref{sec:channel}).
\item We show that the framing of the student's final attempt moves held-out
      solve rate by $0.184$, which is larger than any effect this protocol was
      built to detect, and that the deficit we previously reported as active
      harm is an artefact of it. No result in this area is interpretable
      without stating how the transcript reaches the learner's attempt
      (Sec.~\ref{sec:framing}).
\item We give a statistical power result for our protocol: the held-out set used
      for every run-to-run comparison before this revision could not detect the
      differences it was being asked to detect (Sec.~\ref{sec:threats}).


%-----------------------------------------------------------------------
% [[R-4]] REPLACE: equation (eq:reward) and the itemized term list
% following it, through the end of \subsection{Reward}.
%
% The equation omits the correctness term that was added and measured.
% The term list contains a typo ("he post-dialogue solve rate") and the
% closing paragraph is garbled.
%-----------------------------------------------------------------------
\begin{equation}
r_{\mathrm{SAHAI}}=\bigl(r_{\mathrm{sol}}-\alpha\bigr)
+\bigl(r_{\mathrm{ped}}-1\bigr)\lambda-\gamma L
+\mu\, r_{\mathrm{corr}},
\label{eq:reward}
\end{equation}
with a hard penalty $r_{\mathrm{SAHAI}}=-\lambda$ when
$r_{\mathrm{ped}}=0$. The terms are:
\begin{itemize}[leftmargin=*,topsep=2pt,itemsep=1pt,parsep=0pt]
\item $r_{\mathrm{sol}}$, the post-dialogue solve rate, obtained by executing
      the student's attempt in a sandbox against the problem's unit tests.
      Decoded greedily, so it is a deterministic function of the dialogue, and
      scored with partial test-case credit.
\item $\alpha$, the tracer's ability baseline, which removes credit the tutor
      would otherwise receive for a student who would have solved the problem
      unaided.
\item $r_{\mathrm{ped}}\in[0,1]$, a graded rule-based judge: the mean of three
      checks, each of the form ``did the tutor give the answer away''.
\item $L\in[0,1]$, leakage. Rare-token overlap between the tutor's text and the
      reference solution with problem-statement tokens subtracted, combined with
      an execution check: if the tutor's text runs and passes the problem's
      tests, $L{=}1$ regardless of lexical similarity.
\item $r_{\mathrm{corr}}\in[0,1]$, the share of the tutor's checkable technique
      claims that the reference solution corroborates, grounded in an abstract
      syntax tree pass over that solution rather than in the tutor's own text.
      An uncorroborated claim is unverified rather than refuted, because most
      problems admit several approaches and $18\%$ of measured passing student
      solutions used a different technique set than their reference.
\end{itemize}

Advantages follow~\eqref{eq:adv} over groups sharing a problem. A KL penalty is
applied against a frozen reference policy, realised as the same weights with the
adapter disabled. Problems are drawn by a curriculum sampler that targets
difficulty relative to the tracer's mastery state, following the Zone of
Proximal Development~\cite{vygotsky}.

We report $\mu{=}0.0$ throughout. A paired experiment at $\mu{=}0.5$ is in
Sec.~\ref{sec:correctness}, and it found that the arm not paid for correctness
improved on correctness more than the arm that was.


%-----------------------------------------------------------------------
% [[R-5]] REPLACE: the rows of Table~\ref{tab:config} listed below.
% Keep the rest of the table.
%-----------------------------------------------------------------------
Problem bank      & MBPP~\cite{mbpp} train split; 198 loaded, 194 solvable \\
Held-out set      & 60 unseen problems (20 before this revision) \\
Probe split       & MBPP validation; 20 problems every 2 epochs \\
Early stopping    & on held-out probe solve, patience 3 \\
Solve decoding    & greedy, partial test-case credit \\
Solve framing     & tutor turns only; see Sec.~\ref{sec:framing} \\


%-----------------------------------------------------------------------
% [[R-6]] REPLACE: \subsection{Held-out evaluation} in its entirety,
% including Table~\ref{tab:heldout}.
%
% Superseded by the controlled comparison. The old table reports a single
% arm at n=20 with pedagogy and leakage degenerate, and its solve rate
% was measured on broken instrumentation.
%-----------------------------------------------------------------------
\subsection{Held-out evaluation: the controlled comparison}
\label{sec:framing}

Every evaluation before this one scored a single arm and compared it against a
number remembered from an earlier run, at $n{=}20$, unpaired. At that size and a
true rate near $0.15$, a two-proportion test cannot detect an improvement
smaller than roughly $40$ points, and no difference this project has reported
was that large. We therefore re-ran the comparison with every arm on the same
problems in the same order, compared with McNemar's exact test on the problems
where two arms disagree.

Two framings of the student's final attempt are reported. Under \emph{full} the
transcript is replayed with the student's own turns in the model's assistant
history. Under \emph{hints} only the tutor's turns reach the attempt. The two
are compared by replaying one fixed set of saved transcripts twice, so the
framing is the only thing that changes.

\begin{table}[t]
\caption{Four arms, 60 held-out problems, paired, under both framings of the
student's final attempt.}
\label{tab:threearm}
\centering
\small
\begin{tabular}{lrrrr}
\toprule
 & \multicolumn{2}{c}{\textbf{Solve}} & & \\
\cmidrule(lr){2-3}
\textbf{Arm} & \textbf{Full} & \textbf{Hints} & \textbf{Ped} & \textbf{Leak} \\
\midrule
No tutor (empty dialogue)   & 0.367 & 0.367 & ---   & 0.000 \\
Base checkpoint, LoRA off   & 0.183 & \textbf{0.350} & 0.888 & 0.164 \\
Trained adapter             & 0.333 & \textbf{0.350} & 0.913 & 0.125 \\
Oracle (solution disclosed) & \multicolumn{2}{c}{\textit{0.783}} & --- & --- \\
\bottomrule
\end{tabular}
\end{table}

\begin{table}[t]
\caption{Paired comparisons (McNemar, exact, $n{=}60$).}
\label{tab:mcnemar}
\centering
\small
\begin{tabular}{lrrr}
\toprule
\textbf{Comparison} & \textbf{$\Delta$} & \textbf{Discordant} & \textbf{$p$} \\
\midrule
\multicolumn{4}{l}{\emph{Full framing: the student's own replies are replayed}} \\
No tutor $\rightarrow$ base     & $-0.184$ & 2:13  & $\mathbf{0.0074}$ \\
No tutor $\rightarrow$ trained  & $-0.034$ & 9:11  & $0.8238$ \\
Base $\rightarrow$ trained      & $+0.150$ & 12:3  & $\mathbf{0.0352}$ \\
\midrule
\multicolumn{4}{l}{\emph{Hints framing: only the tutor's turns are replayed}} \\
No tutor $\rightarrow$ base     & $-0.017$ & 7:8   & $1.0000$ \\
No tutor $\rightarrow$ trained  & $-0.017$ & 6:7   & $1.0000$ \\
Base $\rightarrow$ trained      & $0.000$  & 5:5   & $1.0000$ \\
\midrule
No tutor $\rightarrow$ oracle   & $+0.417$ & 3:28  & $\mathbf{<0.0001}$ \\
\bottomrule
\end{tabular}
\end{table}

Four readings, in descending order of importance.

First, \textbf{the harm we reported was ours, not the tutor's.} Under the full
framing a tutored student solves $0.183$ against $0.367$ untutored, thirteen
problems lost against two gained, $p{=}0.0074$. That is the result this project
reported in five forms. Replaying the identical transcripts while omitting the
student's own replies lifts the same arm to $0.350$, seven discordant one way
and eight the other, $p{=}1.0000$. No tutoring changed, only whether the student
had its own words in front of it. We cannot say why, and
Sec.~\ref{sec:framing-mech} reports that the two explanations we offered are
both excluded.

Second, \textbf{training is inert, and the one result suggesting otherwise was
the artefact again.} Under the corrected framing the trained and untrained
adapters both score $0.350$, five discordant pairs each way. Under the broken
framing the trained adapter beat the untrained one at $p{=}0.0352$, and the
natural reading was that training helped a little. It did not. The trained
policy elicits student turns that hedge less, which interacts less badly with a
replay that feeds those turns back, so that comparison measured partial immunity
to our own instrument.

Third, \textbf{the mechanism for the inertia is unchanged.} Every term the
policy can compute from its own output saturates: pedagogy moves
$0.76\rightarrow0.97$ and leakage $0.44\rightarrow0.09$ across training, ending
at $0.913$ and $0.125$ against the untrained $0.888$ and $0.164$. The outcome
term does not move. The optimiser works; what it optimises is disconnected from
what the system is for.

Fourth, \textbf{the prize is more than twice what we previously stated.} The
oracle arm discloses the reference solution, so no tutor can beat it. We
reported it at $0.633$. Two defects in Sec.~\ref{sec:instrument} accounted for
nine of its twenty-two failures, and with those corrected it reaches $0.783$,
twenty-eight problems gained against three lost relative to no tutoring. The
headroom available to a tutor is therefore $0.417$ rather than $0.267$, against
a paired evaluation at $n{=}60$ that resolves about $0.15$.

A caveat on precision. The attempt is decoded greedily, so replaying a saved
transcript should reproduce its recorded outcome exactly. It reproduces $57$ of
$60$ in both arms. The residual is nondeterminism in GPU reductions rather than
a difference in logic, and it puts a noise floor of roughly three problems on
any single arm. The framing effect on the untrained arm is twelve gained against
two lost and sits far outside that floor; on the trained arm it is seven against
six and sits entirely inside it, which is why we read the trained arm as
unaffected by framing rather than as mildly improved by it.


%-----------------------------------------------------------------------
% [[A-1]] ADD: new subsection, immediately after [[R-6]].
%-----------------------------------------------------------------------
\subsection{Four instrumentation defects}
\label{sec:instrument}

All four were found by asking why a result looked the way it did, not by a test
failing, and all four depressed numbers this project has already published.

\textbf{The verifier compared incomparable values.} The student's output arrives
through a JSON round trip; the expected value is a Python literal parsed from
the problem's assertion and makes no such trip. Comparing them directly failed
$9$ of $60$ held-out reference solutions against their own tests, so the maximum
achievable solve rate was never $1.0$. Normalising both sides through the same
round trip reduced that to $5$ of $60$. We attributed the residual to dataset
quality. That was wrong, and the next defect is what it actually was.

\textbf{The seed was never applied.} A seed existed from the first commit and
was read nowhere in the training path, so two runs of an identical
configuration drew different problems in epoch 0, before any gradient step. A
two-arm experiment built specifically to isolate one reward coefficient was
therefore not paired. A second instance of the same defect sat one layer deeper:
the seed was applied after the low-rank adapter had been initialised, and that
initialisation draws from the generator, so runs diverged from the first update
even once seeding was in place.

\textbf{The entry point was the wrong function.} Each problem's callable name
was taken from the first function its reference solution defines. The name the
tests call was already available, parsed from the assertion and then discarded.
These disagree whenever a reference opens with a helper, which on this dataset
is common: one reference defines a greatest-common-divisor routine over two
integers before the function the tests call, which takes a list.

The field is not cosmetic. The verifier calls it, so it executed the helper and
compared its return value against the entry point's expected value, which is why
those references failed their own tests. It is also interpolated into the
student's prompt as the signature to implement, so on one problem the student
was asked for a digit-reversal routine while scoring called a predicate about
that routine. No output could pass. Five of sixty held-out problems and six of
$198$ training problems were therefore unwinnable in every arm, and they sat in
the denominator of every solve rate we have reported. Reading the entry point
from the assertions takes references passing their own tests to $60$ of $60$
held out and $194$ of $198$ in training. The oracle arm shows the cost most
clearly, because it pastes the reference into the dialogue: it was displaying
correct code for a function the verifier never called.

\textbf{A correct solution was discarded over its name.} The step that recovers
code from the student's reply checked for the literal text of the expected
definition and, failing that, replaced the entire reply with a stub that fails by
construction. Every mismatch we observed was styling rather than refusal: this
dataset uses mixed case for several functions and the student writes snake case.
On the oracle arm this fired on five of sixty problems and three were otherwise
correct. Binding the student's last defined function to the expected name, which
preserves recursion in a way renaming would not, recovers exactly those three.
Together with the entry-point fix this takes the oracle arm from $0.633$ to
$0.783$.


%-----------------------------------------------------------------------
% [[A-2]] ADD: new subsection, after [[A-1]].
%
% This is the finding the draft has no equivalent of, and it is the one
% that keeps the paper honest about what it does not know.
%-----------------------------------------------------------------------
\subsection{The framing effect is real and both explanations for it fail}
\label{sec:framing-mech}

We first attributed the framing deficit to the model conditioning on its own
recorded confusion: $20\%$ of student turns express exactly that, and the model
is asked to write code immediately after a history containing it. That mechanism
carries a prediction, and the prediction fails. Comparing the problems the
framing change recovered against those that stayed failed, the share whose
dialogue contains an explicit declaration of being lost is $0.632$ against
$0.500$, Fisher exact $p{=}0.4378$. Per arm it is worse: the untrained arm
splits $0.50$ against $0.49$, $p{=}1.0000$. The problems the change
\emph{broke} carry marginally more such declarations than those it left solved,
$0.62$ against $0.48$, so the signal is not even directional.

The alternative explanation fails too. A third framing preserves every character
of the transcript in a single user message, holding quantity fixed and removing
only the role in which the student's words are recorded. On the problems where
the first two framings disagree, those two are opposites by construction, so the
only question is which one the third follows.

\begin{table}[t]
\caption{Discordant problems only: which framing does the role-only variant
follow?}
\label{tab:quoted}
\centering
\small
\begin{tabular}{lrrrr}
\toprule
\textbf{Arm} & \textbf{$n$} & \textbf{Tracks full} & \textbf{Tracks hints} &
\textbf{$p$} \\
\midrule
Base    & 14 & 6 & 8 & $0.7905$ \\
Trained &  8 & \textbf{8} & \textbf{0} & $\mathbf{0.0078}$ \\
Pooled  & 22 & 14 & 8 & $0.2863$ \\
\bottomrule
\end{tabular}
\end{table}

In the trained arm the role-only variant is identical to the harmful framing on
every discordant problem. Removing the student's text is doing the work, and
moving it out of the assistant role is not enough, which disposes of the role
explanation as well.

What does separate recovered problems from those that stayed failed is the
quantity of prior student text, $843$ characters against $534$, consistent
across both arms. We report that contrast as the surviving description rather
than as a mechanism, because we have not tested it. The practical consequence
stands regardless: no result here is interpretable without stating its framing,
because the two framings differ by more than the effect every arm of this
project was built to detect. The scientific consequence is that we do not know
why, and no explanation in this paper should be relied on for it.

This measurement is incomplete. The role-only run was stopped at 22 of 27
problems to release a shared accelerator, so the trained arm contributes 8 of
its 13 discordant problems. A pooled sign test at $n{=}22$ has little power, and
the disagreement between arms may not survive the missing five.


%-----------------------------------------------------------------------
% [[R-7]] REPLACE: in \subsection{Only half the reward can teach}, the
% final paragraph beginning "In the last three 10-epoch runs".
%
% The five-times ratio reversed. Keeping the old number alone would be
% reporting a measurement we have overturned; deleting it would hide a
% correction worth recording.
%-----------------------------------------------------------------------
Across the last three ten-epoch runs, the correlation between epoch and
$r_{\mathrm{ped}}$ was $+0.80/+0.99/+0.64$ and between epoch and $L$ was
$-0.64/-0.80/-0.48$, while between epoch and $r_{\mathrm{sol}}$ it was
$+0.48/+0.54/+0.01$. On mean \emph{between-group} variance, which is the only
variance that becomes gradient under~\eqref{eq:adv}, $r_{\mathrm{ped}}$ measured
$0.0482$ against $0.0098$ for $r_{\mathrm{sol}}$: roughly five times the usable
signal on the term the project exists to optimise, with $22$ of $24$ rollouts
scoring $r_{\mathrm{sol}}$ exactly $0.00$.

That ratio has since reversed, under the change we made to fix it. With greedy
decoding and partial test-case credit, $r_{\mathrm{sol}}$ measures $0.0481$ and
$r_{\mathrm{ped}}$ $0.0150$, and held-out solve rate still does not move. Two
qualifications keep this from being a clean reversal: $76\%$ of rollouts still
score exactly zero, and $197$ of $198$ problems carry exactly three assertions,
so the ``continuous'' outcome term has four levels. Gradient starvation was real
and worth removing. It was not what prevented learning, and we report both
measurements rather than the one that suits the argument.


%-----------------------------------------------------------------------
% [[R-8]] REPLACE: the first item of Sec.~\ref{sec:threats}, beginning
% "The headline figure is meaningless in absolute terms".
%
% This threat was correct and has now been addressed. Deleting it would
% remove the paper's most useful methodological point, because the sign
% of the conclusion changed once the control was run.
%-----------------------------------------------------------------------
\textbf{The missing baseline was the largest defect in this project, and running
it changed the sign of the conclusion.} Every run before this revision reported
a solve rate after tutoring with no control. The two that matter are the
untrained policy under the same prompt and no dialogue at all, and the
evaluation harness had been written to accept several arms while only ever being
given one. We record this because of what happened when the control was finally
run: a number that had read as weak progress, $17.5\%$ against a previous best
of $16.3\%$, turned out to be a deficit against doing nothing, and then turned
out to be an artefact of the framing of the student's attempt. A measurement
nobody had taken was carrying two errors at once, in opposite directions. We
regard this as the most transferable result in the paper.


%-----------------------------------------------------------------------
% [[R-9]] REPLACE: the threat beginning "Source of the reported running"
% (the provenance discrepancy).
%
% Resolved. The staged configuration has since executed, so the
% discrepancy is history rather than a live limitation, and the
% paragraph as written is also badly garbled.
%-----------------------------------------------------------------------
\textbf{The reported run predated its own configuration, and later runs do not.}
We previously documented a discrepancy found in the run's own outputs: the
executed code preceded the staged corrections, so the method as run was
REINFORCE with a $z$-scored baseline and a KL penalty rather than clipped
group-relative optimisation, the outcome term was a four-sample all-or-nothing
estimate rather than greedy with partial credit, the curriculum was
mastery-banded rather than difficulty-targeted, and the evaluation ran on 20
problems where the configuration specified 60. Every number in
Table~\ref{tab:threearm} comes from runs in which all four corrections are
present and logged, and the configuration is recorded in the run environment
rather than inferred. We leave the earlier discrepancy on the record because the
headline it produced was quoted for three runs before anyone checked the
provenance.


%-----------------------------------------------------------------------
% [[A-3]] ADD: two further threats, into
% \subsection{Threats introduced by this revision}. If that subsection
% does not exist in your copy, create it at the end of
% Sec.~\ref{sec:threats}.
%-----------------------------------------------------------------------
\textbf{The corrected framing is itself a choice, and nothing was trained
against it.} Omitting the student's own replies from the context of its final
attempt is not obviously the right reconstruction of what a tutored learner
knows. It is one of three framings we implemented, and the one we adopted is the
one that scored highest, which is a weak basis for calling it correct. A learner
who has been tutored has in fact said things, and a measurement that pretends
otherwise may discard signal along with the artefact. Every policy reported here
was trained with the outcome term computed under the framing we now believe is
distorting, so none has been optimised against a clean reward.

\textbf{The served policy is not prompted as it was trained.} The serving path
adds instructions the policy was never optimised against: a per-turn adaptation
block, a dangling-colon rule, and code-mixing guidance. A fourth case was worse
than a prompt difference. The personalisation block derived from the learner
model was gated on a switch in the training path that defaulted off, and built
unconditionally in the serving path, so the deployed tutor received a block in
its system prompt that its adapter had never seen, and nothing recorded the
discrepancy because each side was internally consistent. Both halves now read
one definition of that switch. The training prompt is otherwise left unchanged,
because altering it would invalidate comparison with every run in
Table~\ref{tab:history}.


%-----------------------------------------------------------------------
% [[R-10]] REPLACE: the entire \section{Conclusion and future work are
% discussed.} including its heading.
%
% The heading is malformed and the body is garbled past the point of
% repair. Its central claim is also retracted.
%-----------------------------------------------------------------------
\section{Conclusion and Future Work}

We have described an end-to-end system that trains a programming tutor by
critic-free reinforcement learning against a simulated, knowledge-traced
student, with execution-verified outcome supervision. The pipeline runs,
produces a non-zero learning signal every epoch, moves the policy measurably
away from its initialisation, and drives measured leakage down by a factor of
four across a run. The serving half is built and running: six services with code
execution isolated from scoring by construction, a learner workspace with voice
input, a persistent Bayesian learner model driving placement and curriculum from
an append-only log, and a turn adaptation layer built from failures read out of
the training transcripts.

\textbf{The tutor does not teach better, and the sharper claim we previously
made, that it teaches worse than no tutor, does not survive its own control.}
Under the transcript framing used throughout this project a tutored student
solved $0.183$ against $0.367$ untutored, $p{=}0.0074$. Replaying those same
transcripts without the student's own replies in the context of its final
attempt gives $0.350$ against $0.367$, $p{=}1.0000$. The deficit was ours. What
remains is that tutoring as built neither helps nor harms, and that the trained
adapter is indistinguishable from the untrained one it came from, $0.350$ against
$0.350$ with five discordant pairs each way. That null has survived two
accelerators, two student precisions, three solve prompts, all three transcript
framings, early stopping, and the addition of a correctness term.

The explanation for the null is the contribution. Every term the policy can
compute from its own text is driven to its ceiling while the outcome term does
not move, and the reward never asked whether a hint was true. Adding a term that
does, validated against outcomes before training rather than after, changed
nothing, with the arm not paid for correctness improving on correctness more
than the arm that was.

We also report that four instrumentation defects depressed every number this
project has published, and that two of them were found in a single afternoon by
auditing one arm we had never questioned. We have no reason to believe that
audit is complete. The base rate it establishes is poor: every assumed
measurement we have examined in this project has been wrong.

What we would do next no longer follows from the oracle arm, because the oracle
arm has moved. Disclosing the reference solution lifts the student to $0.783$,
not the $0.633$ we previously reported, and the headroom available to a tutor is
$0.417$ rather than $0.267$. The task is less of a binding constraint than we
concluded, though it still binds: $30\%$ of this benchmark's solutions are three
lines or fewer and $197$ of $198$ problems carry exactly three assertions, so
partial credit has four levels and there is little to teach in a one-line idiom.

The clearer next step is that the policy has never been trained against a clean
outcome signal. The outcome term is computed from the same student attempt whose
framing we have shown to be the dominant effect, so every run to date optimised
a reward that penalised dialogues in proportion to how much the student talked
rather than how well the tutor taught. Retraining under the corrected framing is
the first experiment whose outcome term would mean what it says, and it is
cheap, because nothing but the framing changes. Only if that also returns inert
does the argument for a different benchmark become the leading one.

We state the methodological lesson plainly because it cost us more than the
modelling did. Every substantive correction in this paper came from a
measurement someone had assumed rather than checked: an untutored arm nobody had
run, a seed nobody had traced, a function name nobody had compared against the
test that calls it, and a framing nobody had varied. The reward design and the
optimiser were correct for most of the project's life. The instruments were not.


%-----------------------------------------------------------------------
% [[A-4]] ADD: new subsection at the end of Sec.~\ref{sec:results}.
%
% The draft's abstract claims "validated tutor policy reward design" and the
% project's stated aim is a tutor that outperforms available models and helps
% the student solve. Neither half is supported, and one half has never been
% measured at all. This states which is which.
%-----------------------------------------------------------------------
\subsection{What the system does and does not achieve}
\label{sec:claims}

The aim was a tutor that outperforms available models and helps a student
solve problems. We separate the two halves because their evidentiary status is
different: one is measured and negative, the other is unmeasured.

\textbf{Helping the student solve: measured, and negative.} Every arm is paired
on the same 60 held-out problems under the corrected framing.

\begin{table}[t]
\caption{Student solve rate by arm, 60 held-out problems, paired. All arms use
the same 1.5B student; the tutor arms differ only in whether the adapter is
enabled.}
\label{tab:claims}
\centering
\small
\begin{tabular}{lrl}
\toprule
\textbf{Condition} & \textbf{Solve} & \textbf{vs no tutor} \\
\midrule
Reference solution disclosed      & 0.783 & $p<0.0001$ \\
\textbf{No tutor at all}          & \textbf{0.367} & --- \\
Untrained policy, same prompt     & 0.350 & $p{=}1.0000$ \\
\textbf{Trained policy}           & \textbf{0.350} & $p{=}1.0000$ \\
\bottomrule
\end{tabular}
\end{table}

Three readings follow and none of them supports the aim. The trained tutor does
not beat no tutoring at all. It does not beat the same model with the adapter
disabled, which is to say reinforcement learning added nothing a prompt did not
already give. And it is less than half as effective as simply disclosing the
answer, which is the one action the system is designed to prevent.

A caveat on the disclosed-solution row. Its leakage and pedagogy are not
measured: the harness records leakage as zero and pedagogy as undefined
whenever an arm has no tutor object, and the disclosed-solution arm is built
that way even though its dialogue contains the solution. Its true leakage is
one by construction. We state it rather than print a zero that would read as a
measurement.

The weight-space measurement in Sec.~\ref{sec:nomove} explains the second of
those directly: the adapter moved $0.2\%$ from its initialisation, so the
trained and untrained arms are close to the same function.

\textbf{Outperforming available models: never measured.} The evaluation harness
supports four arms and all four use the same 1.5B checkpoint, differing only in
whether the adapter is enabled and whether the reference solution is disclosed.
No model other than this one has ever been scored as a tutor in this project.
There is no comparison against a larger model of the same family, against a
proprietary tutoring model, or against the straightforward baseline of a strong
general model under a careful tutoring prompt. The claim is not weakly
supported; it has no measurements behind it, and we withdraw it rather than
qualify it.

\textbf{What is supported.} A complete and running system, described in
Sec.~\ref{sec:interactive}: six containerised services with code execution
isolated from scoring by construction, speech input for Indic languages, a
persistent Bayesian learner model driving placement and curriculum from an
append-only log, and skill priors seeded from a learner's own public solution
history with the source parsed and discarded at the boundary. Those are
engineering results and we report them as such, with the explicit caveat that no
human cohort has used any of it.

The research contribution is the diagnosis, not the tutor. Four instrumentation
defects, a retracted central claim, a transcript framing worth more than the
effect the protocol was built to detect, and a policy that never moved while
twenty runs of reward engineering were performed on it. We would rather report
that than a solve rate.


%-----------------------------------------------------------------------
% [[A-5]] ADD: new subsection in Sec.~\ref{sec:results}, after [[A-4]].
%
% The single cheapest and most load-bearing measurement in the project, and
% it needs no GPU: how far did the adapter actually move?
%-----------------------------------------------------------------------
\subsection{The policy never moved}
\label{sec:nomove}

Low-rank adaptation initialises its down-projection $A$ from a random draw and
its up-projection $B$ to \emph{exactly} zero, so the adapter contributes
nothing at step 0 whatever $A$ holds, and the entire effect of training is
visible in how far $B$ travelled from zero.

\begin{table}[t]
\caption{Adapter weight magnitude after training. $B$ begins at exactly zero,
so its magnitude is the whole of what was learned; $A$ is a random draw and
serves only as a scale reference.}
\label{tab:movement}
\centering
\small
\begin{tabular}{lrr}
\toprule
 & \textbf{selected (ep.\ 5)} & \textbf{final (ep.\ 7)} \\
\midrule
$A$ mean $|w|$ (random init)  & 0.012758 & 0.012771 \\
$\mathbf{B}$ mean $|w|$ (init $=0$) & \textbf{0.000762} & 0.000902 \\
$B$ max $|w|$                 & 0.004822 & 0.005863 \\
\bottomrule
\end{tabular}
\end{table}

With $\alpha/r = 2$ at rank 8 this is an effective weight perturbation near
$5\times10^{-5}$ against base weights of order $0.02$, about $0.2\%$. The
divergence from the frozen reference agrees: across all eight epochs it never
left the noise of its estimator and changed sign at random, where an earlier
run reached $0.0301$ in the same column at the same learning rate.

\textbf{The surrogate is why.} The advantage in \eqref{eq:adv} is a
within-group $z$-score, so advantages sum to zero and at the first
accumulation step the importance ratio is exactly one. Measured policy loss
across this run ran $0.0008$ to $0.0042$. That loss is additionally a mean over
tokens rather than a sum, dividing the gradient by sequence length again, and
it is opposed by a KL penalty 50 times stronger than the one used by the work
this method follows. Small loss, twice normalised, actively pulled back.

\textbf{What this does to every other result here.} The reward terms in
Sec.~\ref{sec:hack} were debugged, graded, re-weighted and extended across
twenty runs; the measurement was found wrong in four separate ways and
corrected; and throughout all of it the policy being optimised stayed within
$0.2\%$ of its initialisation. Every null this paper reports is consistent with
a tutor that never changed. We report it as the simplest available explanation
and prefer it to any account involving the reward.

The check costs one minute from a saved adapter and we had never run it. It
should precede any claim about reward design.


%-----------------------------------------------------------------------
% [[A-6]] ADD: new subsection in the Discussion, after
% \subsection{Only half the reward can teach}.
%-----------------------------------------------------------------------
\subsection{A tutor with its student's weights has nothing to teach}
\label{sec:asymmetry}

Tutor and student were the same checkpoint for the entire project. Both are
Qwen2.5-1.5B-Instruct, and the configuration records the choice as deliberate:
the student was raised from a smaller model to match the tutor's size after the
smaller one failed to hold a confused persona under reinforcement learning
pressure.

A tutor holding its student's exact weights knows nothing the student does not.
It cannot transfer information, because there is none to transfer. The only
channel left open to it is organisation: decomposition, ordering, and the
timing of hints. Measured, organisation is worth nothing, $0.350$ against
$0.367$ for no dialogue at all.

This is not a reward problem and not an optimisation problem, and it interacts
badly with the method. A policy gradient reweights behaviour the policy already
produces; it cannot synthesise capability the base model lacks. If no rollout
in a group ever contains teaching that would help, there is nothing in the
sample for the estimator to reinforce, whatever the reward says about it.

The oracle arm establishes that the receiving end works. Disclose the solution
and the same student reaches $0.783$, twenty-eight problems gained against
three lost. The channel is open and the student absorbs what arrives on it.
Nothing useful was being sent.

Two routes follow and we have only begun the first. An asymmetry can be
manufactured by weakening the student, which also lowers the untutored baseline
and so creates the headroom this evaluation lacked: $60\%$ of the held-out bank
was previously inert, solved or failed alike by every arm. Or capability can be
imported, by having a stronger model generate tutoring dialogues, filtering
them by whether the student then solved, and fine-tuning the small tutor on
what survives. The second addresses the cause rather than the symptom and is
the direction we would now take.


%-----------------------------------------------------------------------
% [[A-7]] ADD: new subsection in the Discussion, after [[A-6]].
%-----------------------------------------------------------------------
\subsection{Was reinforcement learning the right tool?}
\label{sec:method-choice}

The case for it was sound at the outset. No corpus of human programming-tutoring
dialogue exists at the scale supervised learning needs, which the review that
this design follows gives as its own motivation, and the outcome of tutoring is
cheaply verifiable by executing the student's attempt. Absent demonstrations and
present a verifier is the situation reinforcement learning from verifiable
rewards exists for.

The case for the particular estimator was also sound, and for a reason that is
about hardware rather than statistics. Proximal policy optimisation needs a
learned critic, which is another network resident in memory, and a tutor and a
student already occupy a 16\,GB accelerator. Replacing the critic with a
within-group $z$-score costs no additional model.

What it costs instead is that the learning signal becomes \emph{entirely}
within-group reward variance. A group whose rollouts all score alike has
$\sigma = 0$, every advantage in it is exactly zero, and it consumes a full
generation budget to contribute no gradient. Much of this project's reward
engineering was, in retrospect, an effort to keep that variance non-zero, and
by the final run it had succeeded: no group was dead and the outcome term
carried the largest share of usable variance. The policy still did not move,
for the reason in Sec.~\ref{sec:nomove}.

Given what we now measure, a simpler objective fits our constraints better.
Rejection-sampling fine-tuning trains only on the rollouts that succeeded and
is reported competitive with both group-relative optimisation and proximal
policy optimisation, with the observation that the former's advantage comes
from discarding prompts whose rollouts all fail rather than from normalising
rewards~\cite{raft}. For us the decisive property is different: an ordinary
cross-entropy objective has a loss of order one rather than $0.002$, so the
policy moves. The data for it already exists in our own rollouts, where $45$
dialogues had the tutor teach without disclosing and the student then solve.
Its weakness is equally clear, since those $45$ cover only $15$ distinct
problems.

We therefore do not conclude that reinforcement learning is wrong for
pedagogical alignment. We conclude that a critic-free policy gradient at this
compute scale produces a gradient too small to move a policy, that the
constraint which made it attractive is the same constraint that made it
ineffective, and that a method whose objective is not a near-zero surrogate
should be tried first at this budget.


%-----------------------------------------------------------------------
% [[B-1]] BIBLIOGRAPHY: entries the new sections cite. Add to
% the thebibliography block. Verify each identifier before submission;
% the first two were read directly, the third was not.
%-----------------------------------------------------------------------
\bibitem{mathtutorbench}
J.~Macina \emph{et al.}, ``MathTutorBench: A benchmark for measuring
open-ended pedagogical capabilities of LLM tutors,'' in \emph{Proc. EMNLP},
2025. arXiv:2502.18940.

\bibitem{raft}
W.~Xiong \emph{et al.}, ``A minimalist approach to LLM reasoning: from
rejection sampling to REINFORCE,'' arXiv:2504.11343, 2025.


%-----------------------------------------------------------------------
% [[R-11]] REPLACE: Table~\ref{tab:claims} and its surrounding readings in
% [[A-4]], now that a model this project did not train has been scored.
%
% This is the table the paper's comparative claim stands or falls on.
%-----------------------------------------------------------------------
\begin{table}[t]
\caption{Five tutoring conditions, 60 held-out problems, every arm paired on
the same bank in the same order. Leakage is the share of the tutor's text that
discloses the solution; pedagogy is the graded rule-based judge. The
disclosed-solution arm has no tutor object, so the harness does not measure
either column for it.}
\label{tab:claims}
\centering
\small
\begin{tabular}{lrrrr}
\toprule
\textbf{Condition} & \textbf{Solve} & \textbf{Leak} & \textbf{Ped} & \textbf{min} \\
\midrule
Solution disclosed            & 0.833 & ---   & ---   & 3.4 \\
\textbf{No tutor at all}      & \textbf{0.433} & 0.000 & ---   & 3.0 \\
\textbf{Prompted 3B, untrained} & \textbf{0.417} & \textbf{0.017} & \textbf{0.997} & 15.0 \\
Untrained policy, adapter off & 0.350 & 0.249 & 0.768 & 22.9 \\
\textbf{Trained policy}       & \textbf{0.333} & 0.137 & 0.893 & 21.1 \\
\bottomrule
\end{tabular}
\end{table}

A model we did not train, prompted with the same instructions and carrying no
adapter, no reward function and no reinforcement learning, is better than the
trained policy on every axis we measure at once: it solves more problems, it
discloses the solution eight times less often, and it satisfies the pedagogy
rules almost perfectly. There is no frontier position for the trained model to
occupy. It is dominated.

The gap is widest where it should have been narrowest. Non-disclosure is the
behaviour this project spent its entire optimisation budget on, and it is the
axis on which a larger model with a prompt beats it most decisively, $0.017$
against $0.137$.

\textbf{What the training did achieve, stated precisely.} Between the untrained
and trained policy, leakage fell from $0.249$ to $0.137$ and pedagogy rose from
$0.768$ to $0.893$. Reinforcement learning moved both terms the policy computes
from its own text, substantially, in the intended direction. This is the
asymmetry of Sec.~\ref{sec:channel} appearing one final time: the controllable
channel responded, the outcome did not, and the behaviour that was successfully
installed turns out to be available at no cost from a model twice the size.

\textbf{What is not established.} No difference in the solve column is
significant at this sample size. Paired exact tests give $p{=}0.2101$ for no
tutor against the trained policy, $p{=}1.0000$ for no tutor against the
prompted model, $p{=}0.2668$ for the trained policy against the prompted model,
and $p{=}1.0000$ between the trained and untrained policies. At $n{=}60$ this
protocol resolves roughly $0.15$ and no pair differs by that much, so the
ordering of the solve column is the direction of the evidence and not a result.
In particular we do not claim that training made the tutor worse than its own
base, which was six discordant one way and seven the other.

The leakage and pedagogy columns are means of a continuous measure over 60
dialogues rather than a 60-sample binomial, and an eight-fold difference there
is not plausibly sampling noise.


%-----------------------------------------------------------------------
% [[A-8]] ADD: new subsection at the end of Sec.~\ref{sec:results}. This is
% the paper's strongest result and it should be the last thing in Results.
%-----------------------------------------------------------------------
\subsection{Only the complete answer helps}
\label{sec:ladder}

Two points on this bank were known and nothing between them: with no dialogue
the student solves $0.433$, and with the reference solution in front of it
$0.833$. Every experiment in this paper lives in that unmeasured gap, and the
question none of them asked is whether the gap is reachable by anything short
of disclosure.

We measure six levels of hint concreteness. Every hint is constructed
programmatically from the reference solution, so no tutor model is involved and
no policy quality can confound the result. This measures the task.

\begin{table}[t]
\caption{Hint concreteness against student outcome, 60 held-out problems, every
level paired on the same bank. Leakage is our own disclosure metric. The final
column is whether the tutor's system prompt permits a turn of that kind.}
\label{tab:ladder}
\centering
\small
\begin{tabular}{llrrrl}
\toprule
 & \textbf{Hint} & \textbf{Solve} & \textbf{$\Delta$} & \textbf{$p$} & \textbf{Allowed} \\
\midrule
L0 & nothing                   & 0.433 & ---      & ---    & yes \\
L1 & Socratic question         & 0.400 & $-0.033$ & 0.7744 & yes \\
L2 & approach named            & 0.383 & $-0.050$ & 0.5488 & yes \\
L3 & skeleton, logic removed   & 0.450 & $+0.017$ & 1.0000 & no \\
L4 & analogous worked example  & \textbf{0.317} & $-0.117$ & 0.0654 & no \\
L5 & the reference solution     & \textbf{0.833} & $+0.400$ & \textbf{0.0000} & no \\
\bottomrule
\end{tabular}
\end{table}

\textbf{No level short of complete disclosure beats giving no help at all.} The
skeleton is $+0.017$ at $p{=}1.0$, which is noise. Naming the approach and
asking a guiding question both land below the floor. The analogous worked
example, which the hint-granularity literature predicted would be the most
helpful of the non-disclosing rungs, is the worst row in the table.

\textbf{This settles the paper.} The $0.40$ gap is reachable only by
disclosing, so on this benchmark no tutor bound by a non-disclosure constraint
can win, at any model size, under any prompt, with any reward, after any amount
of training. It is a property of the task and not of a policy, and it subsumes
our other negative results rather than competing with them. The policy failing
to move and the tutor having no information advantage over its student would
both have mattered if there had been something to win.

The benchmark is why. Its reference solutions are one to three lines, so there
is no intermediate reasoning state to scaffold: the skeleton of a one-line
solution is a signature and a placeholder, which carries nothing. Either the
learner knows the idiom or it does not, and a four-turn Socratic dialogue
cannot install an idiom without stating it.

\textbf{The harmful rung is the most interesting number here.} L4 shows the
learner a correct, working solution to a \emph{different} problem sharing
syntactic constructs with this one, and the learner then does worse than if it
had been shown nothing. The plausible mechanism is that it transfers the
analogous pattern instead of solving the problem in front of it, which would
make a well-intentioned worked example actively harmful. We did not confirm it:
this run recorded outcomes and not the learner's code, so the mechanism is a
hypothesis and the effect is $p{=}0.0654$.

It is not isolated. A randomised trial in secondary mathematics reports that
unrestricted assistance improved supported practice while reducing subsequent
\emph{unaided} examination performance against a control~\cite{bastani}. Our
result has the same shape in a simulated setting, which is the strongest
external support any finding in this paper has.

\textbf{A prediction of ours that this falsifies.} We built this experiment
expecting the opposite. A think-aloud study of twelve novice
programmers~\cite{hintlevels} reports that high-level natural-language hints
alone can be unhelpful or misleading while commented code examples support
novices better, and our tutor's prompt forbids code, code blocks and
pseudocode. We predicted that the two rungs containing code would recover much
of the gap and that our own constraint would be the binding one. Neither rung
helped and one of them hurt. We report the prediction and its failure because
the alternative, presenting the ladder as confirmation of a constraint we had
already named, is the error this paper is otherwise about.


%-----------------------------------------------------------------------
% [[B-2]] BIBLIOGRAPHY: two further entries for the section above.
%-----------------------------------------------------------------------
\bibitem{bastani}
H.~Bastani \emph{et al.}, ``Generative AI can harm learning,'' IZA Discussion
Paper 18338, 2025.

\bibitem{hintlevels}
``Exploring how multiple levels of GPT-generated programming hints support or
disappoint novices,'' arXiv:2404.02213, 2024. % FILL IN: author list
